\section{A family of solutions from the radial equation}
\label{sec:numa-radial-family}

The open zero-density band gives more freedom than the two smooth solutions exhibit. Writing a further potential as $G(r)-\rho$ reduces the equation to choosing a slope $G'$ for which one of the three factors vanishes. Switching the vanishing factor at two suitable levels produces a continuum of $C^{1,1}$ solutions for the same background $\omega$ and density $f$.

Retain the coefficients and correction of the preceding construction. On $-1<r<1$ we have
\[
 x=p+a,\qquad y=p+b,\qquad z=z_0,
 \qquad p'=q,\quad z_0\geq|q|.
\]
Both $p$ and $z_0$ are even, $p>0$ in the interior, and $p'<0$ on $(0,1)$. The strict inequality $z_0>|q|$ holds on a neighborhood of zero. Hence
\begin{equation}\label{eq:numa-half-excess}
 \int_0^1z_0(r)\,\mathrm{d}r>\int_0^1|p'(r)|\,\mathrm{d}r=p(0).
\end{equation}

\begin{proposition}\label{prop:numa-radial-continuum}
The same smooth background and smooth density in \cref{thm:degenerate-hermitian-nonuniqueness} admit uncountably many distinct normalized $C^{1,1}$ solutions of
\[
 \mathrm{MA}_\omega(u)=f\omega_X^3.
\]
They include a directly constructed solution invariant under the block exchange $\iota[A:B]=[B:A]$, if the cutoff defining $\rho$ is chosen even.
\end{proposition}
\Needspace{6\baselineskip}
\begin{proof}
Choose $-1<I<0<J<1$ satisfying
\begin{equation}\label{eq:numa-transition-balance}
 \int_I^J z_0(r)\,\mathrm{d}r=p(I)+p(J).
\end{equation}
We first show there are an interval's worth of choices. For fixed $I$, define
\[
 K(I,J)=\int_I^Jz_0(r)\,\mathrm{d}r-p(I)-p(J).
\]
At $I=0$ we have $K(0,0)=-2p(0)<0$ and $K(0,1)>0$ by \eqref{eq:numa-half-excess}. Continuity gives $\epsilon>0$ such that these opposite signs persist for every $I\in[-\epsilon,0]$. For $0<J<1$,
\[
 \partial_JK(I,J)=z_0(J)-p'(J)>0.
\]
The intermediate value theorem therefore supplies a unique $J=J(I)\in(0,1)$ for every $I\in[-\epsilon,0)$. Shrink $\epsilon$ below one.

For any pair satisfying \eqref{eq:numa-transition-balance}, define the slope on $[-1,1]$ by
\[
 h_{I,J}(r)=
 \begin{cases}
 p(r),&-1\leq r\leq I,\\
 p(I)-\displaystyle\int_I^r z_0(s)\,\mathrm{d}s,&I\leq r\leq J,\\
 -p(r),&J\leq r\leq1,
 \end{cases}
\]
and extend it by zero outside that interval. The balance identity makes it continuous at $J$, and continuity at $I$ is immediate. Each piece has bounded derivative. Thus $h_{I,J}$ is globally Lipschitz; it is flat at $\pm1$. Let $G_{I,J}$ be any primitive. It is $C^{1,1}$ and constant at each end of the real line. Consequently $G_{I,J}(r)$ extends to a global $C^{1,1}$ function on $X$, smooth near $X\setminus U$.

On the middle interval the slope lies between $-p$ and $p$. Indeed,
\[
 (h_{I,J}-p)'=-z_0-p'\leq0,
 \qquad
 (h_{I,J}+p)'=-z_0+p'\leq0.
\]
The first difference starts at zero at $I$; the second ends at zero at $J$. These two inequalities give $h_{I,J}\leq p$ and $h_{I,J}\geq-p$, respectively.

Put $\widetilde u_{I,J}=G_{I,J}(r)-\rho$. Its solution form is, almost everywhere on $U$,
\begin{equation}\label{eq:numa-slope-form}
 \omega+\ddc\widetilde u_{I,J}
 =\omega_0+\ddc G_{I,J}(r)
 =(x+h_{I,J})P+(y-h_{I,J})Q+(z+h_{I,J}')R.
\end{equation}
On $I<r<J$, the first two coefficients are nonnegative because $x=p+a$, $y=p+b$ and $-p\leq h_{I,J}\leq p$. The third is zero. On $-1<r<I$, we have $h_{I,J}=p$ and $b=0$, since $I<0$; the second coefficient is zero and the others are nonnegative, using $z_0+p'\geq0$. On $J<r<1$, we have $h_{I,J}=-p$ and $a=0$; the first coefficient is zero and the others are nonnegative, using $z_0-p'\geq0$. Outside the band, $h_{I,J}=h_{I,J}'=0$, so the form equals $\omega_0$.

There is no additional hypersurface term in \eqref{eq:numa-slope-form}: $G_{I,J}$ is $C^{1,1}$, and its first derivative is continuous at both transition points. Its distributional Hessian is the bounded almost-everywhere Hessian. The preceding coefficient inequalities therefore prove $\omega+\ddc\widetilde u_{I,J}\geq0$ as a current.

By \cref{prop:c11-products}, the bounded product of this $C^{1,1}$ potential agrees with the almost-everywhere wedge of its Hessian.

The rank-one formula \eqref{eq:numa-rank-frame} now shows that the cube of \eqref{eq:numa-slope-form} is zero almost everywhere on the band, and equals $\omega_0^3$ off it. These are exactly the defining values of $f\omega_X^3$ in the preceding construction. Thus $\widetilde u_{I,J}$ solves the equation. Subtracting its supremum gives a normalized solution $u_{I,J}$.

Different $I$ in the interval above give different normalized solutions. If $I_1<I_2<0$ and $I_1<r<I_2$, then the slope belonging to $I_2$ is $p(r)$, whereas
\[
 p(r)-h_{I_1,J(I_1)}(r)
       =\int_{I_1}^r(p'(s)+z_0(s))\,\mathrm{d}s>0,
\]
because $p'>0$ on $(-1,0)$. The two primitives therefore cannot differ by a constant. An interval contains uncountably many parameters, proving the first assertion. None of these slopes is identically $p$ or identically $-p$, so the constructed solutions are also distinct from the original two smooth solutions.

Finally, the function
\[
 L(t)=\int_0^t z_0(r)\,\mathrm{d}r-p(t),\qquad 0\leq t\leq1,
\]
is negative at zero, positive at one, and strictly increasing on $(0,1)$. Its unique zero $t_*$ gives an admissible symmetric pair $I=-t_*$, $J=t_*$. Evenness of $z_0$ and $p$ then makes $h_{I,J}$ odd: on the middle interval it is $-\int_0^r z_0(s)\,\mathrm{d}s$. Choose its primitive even. If $\rho$ is even, $G_{I,J}(r)-\rho$ and its normalized version are fixed by $\iota$. This proves the explicit invariant-solution assertion.
\end{proof}

The transition hypersurfaces lie inside the zero-density band. The resulting family also raises a selection question for positive-exponent approximation. For the symmetric data, \cref{thm:numa-selected-limit} proves that it selects the symmetric member just constructed.
